% ====================================================================
% FILE: content/14_disciplines_extended.tex
% Disziplinübergeordnete Instantiierungen des Kontinuumrahmens
% Kernausgabe 1.3 Flagship
% ====================================================================

\section{Disziplinübergeordnete Instantiierungen des Kontinuumrahmens}
\label{sec:disciplines-extended}

Dieses Kapitel prueft eine Behauptung begrenzter Darstellung: Ausgewaehlte
wissenschaftliche und formal-theoretische Domänen koennen durch das Tupel
\[
  (\Omega, A, P(t), J(t), \Theta, \partial\Omega, C, k(t))
\]
auf eine Weise ausgedrueckt werden, die mit der strukturellen Hierarchie der
Niveaus \(K_0\)--\(K_{12}\), der Schwellenlandchaft, der Operatorfamilie und der
Verzweigungstopologie aus Abschnitt~\ref{sec:branching-topology} kompatibel ist.
Hier steht \(k(t)\) lediglich dafuer da, den Domänen-Tuple lesbar zu halten;
lokal bezeichnet es das entsprechende \(k(K_x,t)\).
Ebenso können spätere Domänenblöcke Symbole wie \(P_x\) oder \(J_x\)
auch ohne \(t\) schreiben, wenn der Absatz die Familie von Potenzialen oder
Flussen beschreibt und nicht eine einzelne Zeitflaeche.
Jede domanenspezifische Zuordnung wird hier unter einem kanonischen Label behandelt:
der \emph{candidate-embedding rule}. Ausdruecke wie bounded
representation test, structural-compatibility screen, bounded embedding pass und
bounded embedding criterion sind Ortsbeschreibungen dieser Regel, keine
separaten Kriterien. Ziel ist es zu zeigen, wie gut untersuchte Klassen von
Systemen als Kontinua innerhalb der strukturellen Axiome von OC dargestellt werden
koennen, nicht zu behaupten, dass jedes Objekt einer Disziplin durch eine
einzige OC-Kodierung ausgenuehrt ist
\parencite{anderson_more_is_different}.

Dies ist eine Strukturverträglichkeitsthese, nicht eine Behauptung ueber
operativen Nutzen. Der praktische Vergleich wird auf Abschnitt
\ref{sec:oc-core-1-3-practical-utility} und Anhang
\ref{sec:oc-core-1-3-practical-utility-atlas} verschoben. Dieses Kapitel fragt,
ob eine Domäne als begrenztes Kontinuum dargestellt werden kann; es behauptet
nicht, dass jede solche Domäne bereits ein geschlossenes Theorem-Paket,
halte-out Replay, Comparator-Audit oder deploymentbereites Verfahren besitzt.

% ====================================================================
\subsection{Prinzipien fuer Domänen-Einbettung}
% ====================================================================

Um eine Disziplin in OC zu repräsentieren, müssen folgende strukturelle
Ausrichtungsbedingungen erfuellt werden.

\paragraph{1. Abbildung von Objekten auf Kontinuumskomponenten.}
Eine Domäne hat eine Kandidatenkontinuumdarstellung, wenn ihre Objekte abgebildet
werden koennen auf:
\begin{itemize}
  \item \(\Omega\): zulaessige Konfigurationen,
  \item \(A\): Achsen, die wesentliche Freiheitsgrade bestimmen,
  \item \(P\): Potenziale, die Kraefte, Energien, Kosten oder semantische
        Spannung kodieren,
  \item \(J\): Flüsse (Materie, Ladung, Wahrscheinlichkeit, Information,
        Ressourcen),
  \item \(\Theta\): Schwellen, die Stabilität, Phasenwechsel und Kollaps
        steuern,
  \item \(\partial\Omega\): Viabilitaetsgrenzen,
  \item \(C\): Zyklen zur Erhaltung von Persistenz,
  \item \(k(t)\): Continuumness (Viabilitätsmass).
\end{itemize}

\paragraph{2. Einbettungsraumrestriktionen.}
Eine gültige Instanziierung erfordert sowohl
\(\Omega(K_x)\subseteq\Omega(M_x)\) als auch \(A(K_x)\subseteq A(M_x)\). Wenn
nicht die erforderlichen Zustände oder Achsen im Einbettungsraum vorhanden sind,
ist die vorgeschlagene Instanziierung strukturell unzulaessig.

\paragraph{3. K-Level-Entsprechung.}
Eine Disziplin ist auf Niveau \(K_x\) nur dann begrenzt eingebettet, wenn:
\begin{itemize}
  \item ihre dominanten Achsen den Achsen von \(K_x\) entsprechen,
  \item ihre Schwellen die von niedrigeren Ebenen verfeinern,
  \item ihre Flüsse und Zyklen zur Operatorfamilie von \(K_x\) passen,
  \item ihr Verhalten die dimensionsbezogene Monotonie währt.
\end{itemize}

\paragraph{4. Strukturelle Kompatibilität.}
Domäneninstanziierung erfordert:
\begin{itemize}
  \item interne Konsistenz mit OC-Schwellen,
  \item Existenz wenigstens eines unterstuetzenden Zyklus
        \(C_{\mathrm{support}}\),
  \item nicht-leere zulaessige Region \(\Omega\neq\emptyset\),
  \item Erweiterbarkeit der Domänendynamik durch die Operatorfamilie.
\end{itemize}

Ein Vorschlag scheitert, wenn irgendeine Pruefung misslingt: die Abbildung hat eine
ausgeschnittene zulaessige Region, erforderliche Achsen fehlen im Einbettungsraum,
angegebene Schwellen während des Verfahrens nicht erfuellbar sind oder die
vorgesehenen Zyklen das Objekt nicht von der relevanten Grenze fernhalten. Diese
sind Ablehnungsbedingungen, nicht bloss kosmetische Omissionsfehler.

Die folgenden Abschnitte zusammenfassen begrenzte Kandidatenrepräsentationen
ueber Disziplinen.

% ====================================================================
\subsection{Physik (Niveau \texorpdfstring{$K_2$}{K2})}
% ====================================================================

Physikalische Systeme instanziieren \(K_2\), wenn Feldkonfigurationen,
Ordnungsparameter und Kohärenzstrukturen die relevanten Achsen bilden.
Die Diskussion ist bewußt auf Standard-Phänomene, Renormierung und
Perkollation beschränkt, nicht auf eine uneingeschriebene Behauptung über die
ganze Physik \parencite{goldenfeld_lectures,wilson_renormalization,stauffer_percolation}.

\paragraph{Zustandsraum und Achsen.}
\begin{itemize}
  \item \(\Omega_{2}\): Feldkonfigurationen (Skalar-, Vektor- und
        Eichfelder).
  \item \(A_{2}\): räumliche Achsen, Modenachsen, interne Symmetrieachsen,
        Ordnungsachsen.
\end{itemize}

\paragraph{Potenziale und Flüsse.}
\begin{itemize}
  \item \(P_2\): physikalische Energiefunktionale (Wirkung, freie Energie).
  \item \(J_2\): diffusive, konvektive und feldtheoretische Ströme.
\end{itemize}

\paragraph{Schwellenwerte.}
\begin{itemize}
  \item \(\Theta_{\mathrm{crit}}\): Phasengrenzen,
  \item \(\Theta_{\mathrm{BKT}}\): Berezinskii-Kosterlitz-Thouless
        Vortex-Entwirrungs-Schwelle,
  \item \(\Theta_{\mathrm{mass}}\): Kohärenzschwelle für Massenerzeugung,
  \item \(\Theta_{\mathrm{death}}\): Konfinenz- oder Entkopplungsbedingungen.
\end{itemize}

\paragraph{Randstrukturen.}
Innerhalb dieser begrenzten Abbildung entspricht \(\partial\Omega\) kritischen
Flächen, Perkolationsgrenzen und Phasenseparationsschnittstellen.

\paragraph{Zyklen.}
Topologische Solitonen, Vortex-Paare, Renormierungzyklen und oszillatorische
Feldmoden sind Instanzen von \(C_2\).

\paragraph{Begrenztes Einbettungskriterium (Physik -> \texorpdfstring{\(K_2\)}{K2}).}
Ein physikalisches System erfährt das begrenzte Einbettungskriterium auf
\(K_2\), wenn:
\begin{itemize}
  \item seine Felder einen zusammenhaengenden Konfigurationsraum \(\Omega_2\)
        definieren,
  \item sein Energiefunktional Schwellen definiert, die mit OC kompatibel sind,
  \item Flüsse Erhaltungssatze via \(J_2\) ausdrücken,
  \item kohärente Zyklen strukturelle Persistenz erhalten.
\end{itemize}

% ====================================================================
\subsection{Chemie und Entstehung des Lebens (Niveaus \texorpdfstring{$K_3$}{K3}--\texorpdfstring{$K_4$}{K4})}
% ====================================================================

Chemische Systeme werden auf \(K_3\) abgebildet, protocyte Systeme auf
\(K_4\).
Diese Abbildungen sind durch den Abschluss von Reaktionsnetzen, katalytische
Organisation und membrankonstraint präbiotische Struktur beschraenkt
\parencite{kauffman_autocatalytic,raf_networks}.

\paragraph{Chemische Kontinua (\texorpdfstring{\(K_3\)}{K3}).}
\begin{itemize}
  \item \(\Omega_{\mathrm{chem}}\): Konzentrationsvektoren, Reaktionszustände,
        katalytische Konfigurationen.
  \item \(A_{\mathrm{chem}}\): Spezieskonzentrationen, Umweltparameter
        (pH, Salinität, Temperatur).
  \item \(P_{\mathrm{chem}}(t)\): chemische Potentiale, Affinitätsgradienten,
        Konzentrationsdisequilibrien.
  \item \(J_{\mathrm{chem}}\): Reaktionsraten, Transportflüsse.
  \item \(\Theta_{\mathrm{act}}\): katalytische Aktivierungsschwelle.
  \item \(\Theta_{\mathrm{closure}}\): reflexiv auto-katalytischer und
        nahrungs-generierter (RAF)-Abschluss-Schwellenwert.
  \item \(C_{\mathrm{chem}}\): Reaktionszyklen, RAF-Abschluss-Schleifen.
\end{itemize}

\paragraph{Protocell-Kontinua (\texorpdfstring{\(K_4\)}{K4}).}
\begin{itemize}
  \item \(\Omega_{4}\): Vesikelzustände, Membrankonfigurationen, interne
        chemische Zusammensetzungen.
  \item \(A_4\): Krümmung, Oberflächenzug, Permeabilität, Ladung.
  \item \(P_4(t)\): osmotische, chemische, elektrochemische und
        Krümmungspotentiale.
  \item \(\Theta_{\mathrm{perm}}, \Theta_{\mathrm{osm}}, \Theta_{\mathrm{curv}}\):
        Membrantieffizienz-Schwellen.
  \item \(J_4\): osmotische Flüsse, Ionenflüsse, metabolische Flüsse.
  \item \(C_4\): metabolische Zyklen, Membranhaltungszyklen,
        gradient-erhaltende Zyklen.
\end{itemize}

\paragraph{Begrenztes Einbettungskriterium (Chemie -> \texorpdfstring{\(K_3\)}{K3}).}
Ein chemisches System erfüllt das begrenzte \(K_3\)-Kriterium, wenn:
\begin{itemize}
  \item Reaktionsnetze eine nichtleere zulaessige Region bilden,
  \item katalytische und Abschluss-Schwellen ohne membranspezifische Achsen
        darstellbar sind,
  \item Zyklen Reaktionsschluss und Konzentrationskontrolle aufrechterhalten,
  \item Einbettungsrestriktionen die chemischen Achsen des Weges stützen.
\end{itemize}

\paragraph{Begrenztes Einbettungskriterium (Protocells -> \texorpdfstring{\(K_4\)}{K4}).}
Ein protocellulares System erfüllt das begrenzte \(K_4\)-Kriterium, wenn:
\begin{itemize}
  \item membranduell gestuetzte chemische Dynamik eine nichtleere zulaessige
        Region bildet,
  \item Membrandrücke \(\partial\Omega_4\) definieren,
  \item Zyklen Konzentrationsgradienten und Integrität der Kompartimente
        aufrechterhalten,
  \item Einbettungsrestriktionen sowohl chemische als auch membranbezogene
        Achsen stützen.
\end{itemize}

% ====================================================================
\subsection{Biologie (Niveau \texorpdfstring{$K_5$}{K5})}
% ====================================================================

Biologische Zellen, über den protocellulären \(K_4\)-Weg hinaus, entstehen dann,
wenn Protocellen Erregbarkeit und strukturierte Regulation erwerben.
Hier richtet sich die Zielsetzung auf die begrenzte Einbettung von excitablem und
morphogenetischem Organisationsbereich, nicht auf eine unbeschränkte Theorie des
Lebens als Ganzes \parencite{turing_morphogenesis,hopfield_nn}.

\paragraph{Zustandsraum und Achsen.}
\begin{itemize}
  \item \(\Omega_5\): Verteilungen von Membranpotentialen, Kanalzuständen,
        Iongradienten.
  \item \(A_5\): Exzitabilitätsachsen (Membranpotential \(\Delta V\), Ionen-
        Konzentrationen, Gating-Variablenachsen).
\end{itemize}

\paragraph{Potenziale und Flüsse.}
\[
  P_5: \text{elektrochemische Potentiale},\qquad
  J_5: \text{Ionenströme, Stoffwechselströme}.
\]

\paragraph{Schwellen.}
\[
  \Theta_{\mathrm{exc}},\quad
  \Theta_{\mathrm{refr}},\quad
  \Theta_{\mathrm{grad}},\quad
  \Theta_{\mathrm{death}}.
\]

\paragraph{Zyklen.}
Zu den relevanten Zyklen gehoeren der Proto-Spike-Zyklus
\(C_{\mathrm{spike}}\), Stoffwechselzyklen und elektrische Oszillations-
schleifen.

\paragraph{Begrenztes Einbettungskriterium (Biologie -> \texorpdfstring{\(K_5\)}{K5}).}
Ein biologisches System erfüllt das begrenzte \(K_5\)-Kriterium, wenn:
\begin{itemize}
  \item Exzitabilitaet nichttriviale \(\Delta V\)-Achsen definiert,
  \item Schwellen Schwellwerte fuer Feuern und Refraktärae definieren,
  \item Zyklen Homöostase und Signalübertragung erhalten,
  \item Einbettung ionische, energetische und membranstuetzende Achsen enthaelt.
\end{itemize}

% ====================================================================
\subsection{Kognition (Niveau \texorpdfstring{$K_6$}{K6})}
% ====================================================================

Die Abschnitte \(K_6\)--\(K_{10}\) unten sind strukturelle
Kandidaten-Embeddings.
Sie sind Ausblicke auf die Grenzrepräsentation, solange
Abschnitt~\ref{sec:oc-core-1-3-practical-utility} oder
Anhang~\ref{sec:oc-core-1-3-practical-utility-atlas} einen engeren operativen
Lane explizit freigibt. Sie sind nicht unmittelbar anwendbare praktische
Behauptungen.

Kognitive Systeme entstehen, wenn Repräsentation, Bindung und Vorhersagemodelle
zu den dominanten Achsen werden.
Die Darstellung orientiert sich an klassischer Arbeit über Organisationsstrukturen
von Verhalten und Vorhersageverarbeitung, nicht an der Behauptung, dass OC die
Kognitionsforschung vollständig ersetzt
\parencite{hebb_organization,friston_free_energy}.

\paragraph{Zustandsraum und Achsen.}
\begin{itemize}
  \item \(\Omega_6\): repräsentationale Konfigurationen, aktivierte
        konzeptionelle Zustände, Vorhersagezustände.
  \item \(A_6\): repräsentatorische Achsen, Bindungsachsen,
        prädiktive Achsen.
\end{itemize}

\paragraph{Potenziale und Flüsse.}
\begin{itemize}
  \item \(P_6\): Vorhersagefehler, semantische Spannung, Konfidenz.
  \item \(J_6\): Aufmerksamkeitsflüsse, Inferenzflüsse, Gedächtnis-
        transitions.
\end{itemize}

\paragraph{Schwellen.}
\[
  \Theta_{\mathrm{pred}},\quad
  \Theta_{\mathrm{bind}},\quad
  \Theta_{\mathrm{coh}},\quad
  \Theta_{\mathrm{expr}}.
\]

\paragraph{Zyklen.}
Zu den relevanten Zyklen gehoeren Aufmerksamkeits-Vorhersage-Zyklen,
Gedächtnis-Konsolidierungszyklen und Modelle der Stabilitaetsschleifen.

\paragraph{Begrenztes Einbettungskriterium (Kognition -> \texorpdfstring{\(K_6\)}{K6}).}
Ein kognitives System erfüllt das begrenzte \(K_6\)-Kriterium, wenn:
\begin{itemize}
  \item es kohärente innere Modelle innerhalb der
        expressivity-Schwellwerte aufrechterhaelt,
  \item Flüsse den Vorhersagefehler reduzieren und konzeptionelle Struktur
        tragen,
  \item Zyklen Aufmerksamkeit, Erinnerung und Aktualisierung regulieren,
  \item Grenzen \(\partial\Omega_6\) Modellgrenzen kodieren.
\end{itemize}

% ====================================================================
\subsection{Gesellschaft (Niveau \texorpdfstring{$K_7$}{K7})}
% ====================================================================

Soziale Systeme instanziieren \(K_7\), wenn Normen, Institutionen und Vertrauen
die dominanten strukturellen Achsen bilden.
Der Zielumfang ist schwellengetriebene Kollektivverhalten, netzwerkte Koordination
und institutionelle Persistenz \parencite{granovetter_threshold,schelling_micromotives,newman_networks}.

\paragraph{Zustandsraum und Achsen.}
\begin{itemize}
  \item \(\Omega_7\): Kommunikationsgraphen, Rollenstrukturen,
        institutionelle Zustände.
  \item \(A_7\): normative Achsen, Vertrauenachsen, Institutionenachsen.
\end{itemize}

\paragraph{Potenziale und Flüsse.}
\[
\begin{aligned}
  P_7 &: \text{Vertrauen, Legitimität, symbolisches Kapital},\\
  J_7 &: \text{Kommunikation, Ressourcentransfer und Normtransmission}.
\end{aligned}
\]

\paragraph{Schwellen.}
\[
  \Theta_{\mathrm{trust}},\quad
  \Theta_{\mathrm{leg}},\quad
  \Theta_{\mathrm{role}}.
\]

\paragraph{Zyklen.}
Zu den relevanten Zyklen gehoeren institutionelle Zyklen, normhafte Zyklen
und Koordinationszyklen.

\paragraph{Begrenztes Einbettungskriterium (Gesellschaft -> \texorpdfstring{\(K_7\)}{K7}).}
Ein soziales System erfüllt das begrenzte \(K_7\)-Kriterium, wenn:
\begin{itemize}
  \item Vertrauen und Normen Viabilitätsschwellen definieren,
  \item Institutionen Zyklen von Governance und Koordination aufrechterhalten,
  \item Kommunikations- und Ressourcenflüsse die Einbettungsrestriktionen
        respektieren,
  \item soziale Grenzen eine nichtleere \(\partial\Omega_7\) bilden.
\end{itemize}

% ====================================================================
\subsection{Zivilisatorische Systeme (Niveau \texorpdfstring{$K_8$}{K8})}
% ====================================================================

Zivilisatorische Systeme erweitern soziale Kontinua auf Infrastrukturen,
Energieflüsse und großskalige Koordination.
Der Zweck ist beschraenkt auf grossskalige organisierte Systeme mit
netzbasierten Infrastrukturen und Stabilitaetsrestriktionen
\parencite{newman_networks,may_stability_complexity}.

\paragraph{Zustandsraum und Achsen.}
\begin{itemize}
  \item \(\Omega_8\): Infrastrukturnetze, wirtschaftliche Zustände,
        Bevölkerungsverteilungen.
  \item \(A_8\): Infrastrukturelle Achsen, Energieflussachsen,
        Kommunikationsachsen, Komplexitätsachsen.
\end{itemize}

\paragraph{Potenziale und Flüsse.}
\[
  P_8: \text{Energieverfugbarkeit, Ressourcengradienten, Risiko-Potentiale},\qquad
  J_8: \text{Handel, Energiefluss, Datenfluetze}.
\]

\paragraph{Schwellen.}
\[
  \Theta_{\mathrm{stress}},\quad
  \Theta_{\mathrm{capacity}},\quad
  \Theta_{\mathrm{cohesion}}.
\]

\paragraph{Zyklen.}
Relevante Zyklen sind wirtschaftliche Zyklen, Infrastruktur-Renewal-Zyklen
und Wissensweitergabe-Schleifen.

\paragraph{Begrenztes Einbettungskriterium (Zivilisation -> \texorpdfstring{\(K_8\)}{K8}).}
Ein Zivilisationssystem erfüllt das begrenzte \(K_8\)-Kriterium, wenn:
\begin{itemize}
  \item Energie- und Strukturverteilungen \(P_8\) definieren,
  \item Infrastrukturfloesse stabile globale Zyklen bilden,
  \item Institutionen auf \(K_7\) konsistent als Unterstrukturen eingebettet sind,
  \item großräumige Grenzen eine globale \(\partial\Omega_8\) definieren.
\end{itemize}

% ====================================================================
\subsection{Theorie und formale Systeme (Niveaus \texorpdfstring{$K_9$}{K9}--\texorpdfstring{$K_{10}$}{K10})}
% ====================================================================

Theorie- und meta-theoretische formale Systemstrukturen belegen die obere Route in
diesem Kapitel. Dieser Unterabschnitt gibt ein OC-internes Schema für
\(K_9\)--\(K_{10}\), mit lokaler Unterstützung durch
Abschnitte~\ref{sec:klevels-full} und~\ref{sec:oc13-foundational-consistency}.
Es wird nicht behauptet, dass ein vollständig abgeschlossenes externes
Embedding-Theorem fuer jede Theorie oder jedes formale System existiert. Rekursive
strukturelle Kontinua bleiben die Ebene \(K_{11}\), ausserhalb der Regel dieses
Kapitels.

\paragraph{Theoretische Kontinua (\texorpdfstring{\(K_9\)}{K9}).}
\begin{itemize}
  \item \(\Omega_9\): Raum von Modellen, Formalismen und Repräsentationen.
  \item \(A_9\): konzeptionelle, formale und Modellierungsachsen.
  \item \(P_9\): Kohärenzpotentiale, empirische Spannung, erklärende
        Kraft.
  \item \(J_9\): Modellrevisionsflüsse, Anomaliebehandlung, Argument-
        transfer flows.
  \item \(\Theta_{\mathrm{expr}}, \Theta_{\mathrm{coh}}\): Schwellen fuer
        Konsistenz und Expressivität.
  \item \(\partial\Omega_9\): Inkonsistenz, Expressivitaetsversagen oder
        nicht behobene Anomaliegrenzen.
  \item \(C_9\): Paradigmenzyklen, Forschungszyklen.
  \item \(k_9(t)\): Kohärenz- und Validierungsbandbreite der Continuumness.
\end{itemize}

\paragraph{Metatheoretische formale Kontinua (\texorpdfstring{\(K_{10}\)}{K10}).}
\begin{itemize}
  \item \(\Omega_{10}\): Raum metatheoretischer Zustände, formale
        Systeme, Beweiszustände, Modelle-von-Modellen-Objekte und
        rekursionsgebundene Sprachen.
  \item \(A_{10}\): Metaachsen, Rekursionsachsen, Konsistenzachsen,
        Expressivitätsachsen und Theorievergleichachsen.
  \item \(P_{10}\): Suchflusskapazität, Konsistenzmarge, formale
        expressive Kapazität und metatheoretisches Kohärenzpotenzial.
  \item \(J_{10}\): Beweissuchflüsse, formale Updateflüsse,
        Metamodelluebersetzungen und Rekursionssteuerungsflüsse.
  \item \(\Theta_{\mathrm{cons}}, \Theta_{\mathrm{rec}}, \Theta_{\mathrm{meta}}\):
        Schwellen fuer Konsistenz, begrenzte Rekursion und Metaebenenkohärenz.
  \item \(\partial\Omega_{10}\): Inkonsistenz, unentscheidbare Überdehnung,
        Rekursionsrandflächen oder Metaebenen-Grenzfehler.
  \item \(C_{10}\): Beweisschleifen, Rekursionszyklen, Formupdate-
        Zyklen und Metamodell-Auditzyklen.
  \item \(k_{10}(t)\): formale und metatheoretische Lebensfähigkeit unter
        Konsistenz-, Rekursions-, Selbstreferenz- und Beweisdurchsatz-
        Restriktionen.
\end{itemize}

\paragraph{Kandidaten-Einbettungsregel (Theorie -> \texorpdfstring{\(K_9\)}{K9},
Metatheoretisches formales System -> \texorpdfstring{\(K_{10}\)}{K10}).}
Theorie- und metatheoretische formale Routen erfuellen die Kandidatenregel,
wenn die relevante \(K_9\)- oder \(K_{10}\)-Struktur folgende
Bedingungen erfuellt:
\begin{itemize}
  \item konzeptionelle und formale Achsen bilden ein valides Achsenset,
  \item Kohärenzbedingungen definieren mit OC kompatible Schwellen,
  \item Zyklen von Argumentation, Modellierung und Revision erhalten
        Viabilität,
  \item Einbettungsraeume stützen die erforderlichen expressiven Freiheitsgrade.
\end{itemize}

% ====================================================================
\subsection*{Zusammenfassung}

Dieses Kapitel formuliert begrenzte disziplinübergreifende Kandidaten-
Einbettungskriterien fuer Physik, Chemie, Protocells, Biologie, Kognition,
Gesellschaft, Zivilisation, Theorie und formale Systeme. Jede begrenzte Route
wird mit demselben strukturellen Tuple auf der Detaillierungsebene des jeweiligen
Unterabschnittes gemappt, aber jede Zuordnung wird nur als Strukturverträglichkeit-
behauptung gesehen, nicht als Beweis, dass OC bereits die bestmoeglichen lokalen
Modelle in jeder Disziplin ersetzt. Diese stärkere operative und Vergleichsfrage
wird spaeter in den expliziten Unterstuetzungsklassen in
Abschnitt~\ref{sec:oc-core-1-3-practical-utility} und der zeilenweisen
Atlasdarstellung in
Anhang~\ref{sec:oc-core-1-3-practical-utility-atlas} behandelt.

% END OF FILE

